\documentclass[11pt]{article}
\usepackage[margin=1in]{geometry}
\usepackage{amsmath,amssymb,amsthm}
\usepackage{array}
\usepackage{parskip}
\usepackage{booktabs}

\newtheorem{theorem}{Theorem}
\newtheorem{lemma}[theorem]{Lemma}
\newtheorem{corollary}[theorem]{Corollary}
\newtheorem{proposition}[theorem]{Proposition}
\theoremstyle{definition}
\newtheorem{definition}[theorem]{Definition}
\theoremstyle{remark}
\newtheorem{remark}[theorem]{Remark}

\newcommand{\F}{\mathbb{F}}
\newcommand{\Pn}{\mathcal{P}}
\newcommand{\isect}[1]{\bigcap_{i\in #1}}

\title{Disproof of conjecture \texttt{00000002175}}
\author{earthking11}
\date{2026-09-13}

\begin{document}
\maketitle

\section*{Verdict: FALSE}

We disprove conjecture \texttt{00000002175} as stated. The refutation is
unconditional and elementary, and is already complete on a four-element ground
set: there is a 3-wise odd-intersecting family of size $4$, whereas the
conjecture asserts the maximum is $2^{n-3} = 2$. Computations of the exact
maxima for $n \le 7$ contradict the claimed formula in both directions, and an
exhaustive check of all affine cosets shows that no affine coset is an extremal
family, so the characterisation clause is false as well.

\section{The conjecture}

\begin{definition}[3-wise odd-intersecting]
A family $\mathcal{F} \subseteq \Pn([n])$ of subsets of an $n$-element ground
set $[n] = \{1, \dots, n\}$ is \emph{3-wise odd-intersecting} if for every three
\emph{distinct} members $A, B, C \in \mathcal{F}$ the intersection $A \cap B
\cap C$ has odd cardinality.
\end{definition}

Quoted from \texttt{conjectures/00000002175.md}:

\begin{quote}
\textbf{Definition:} A 3-wise odd-intersecting family is one where the
intersection of any three members is odd. \textbf{Conjecture:} The maximal size
is always $2^{n-3}$; and the extremal families consist of indicator functions of
3-dimensional affine subspaces. (3-wise odd-intersection affine extremal)
\end{quote}

The original statement is filed in English and Chinese; the Chinese text reads
(the conjecture-file version is reproduced verbatim in \texttt{README.md}).
The size claim (the maximum is always $2^{n-3}$) and the characterisation claim
(the extremal families are the indicator functions of 3-dimensional affine
subspaces) both fail.

\section{The counterexample at $n = 4$}

\begin{theorem}\label{thm:witness}
On the ground set $\{1,2,3,4\}$, the four-member family
\[
F \;=\; \bigl\{\, \{1,2\},\ \{1,3\},\ \{1,4\},\ \{1,2,3,4\}\,\bigr\}
\]
is 3-wise odd-intersecting, yet $|F| = 4 > 2 = 2^{4-3}$.
\end{theorem}

\begin{proof}
There are $\binom{4}{3} = 4$ triples of distinct members, and each meets in
exactly $\{1\}$, of odd size $1$:
\[
\begin{array}{llll}
\{1,2\}\cap\{1,3\}\cap\{1,4\} &= \{1\}, & \{1,2\}\cap\{1,3\}\cap\{1,2,3,4\} &= \{1\},\\[2pt]
\{1,2\}\cap\{1,4\}\cap\{1,2,3,4\} &= \{1\}, & \{1,3\}\cap\{1,4\}\cap\{1,2,3,4\} &= \{1\}.
\end{array}
\]
Indeed, every one of the four members contains $1$ (so every intersection
contains $1$), and no other point lies in all three members of any triple. Hence
$F$ is 3-wise odd-intersecting. Since $|F| = 4$ while $2^{4-3} = 2$, the size
claim of the conjecture fails.
\end{proof}

\begin{corollary}
The true maximum at $n = 4$ is at least $5$, so the conjecture's value
$2^{4-3} = 2$ is too small by at least $3$. A $5$-element witness is
\[
\bigl\{\, \{1\},\ \{1,2\},\ \{1,3\},\ \{1,4\},\ \{1,2,3\}\,\bigr\},
\]
in which every triple of distinct members meets in $\{1\}$ (a triple containing
$\{1\}$ always meets in $\{1\}$; among the remaining four sets, which all
contain $1$, any three have intersection exactly $\{1\}$).
\end{corollary}

The conjecture is thus refuted by the four-member family alone; the five-member
family shows the failure is even larger than the witness indicates.

\section{The true maxima versus $2^{n-3}$}

The exact maximum $M(n)$ was computed in two independent ways (both implemented
in \texttt{reproduce.py}):

\begin{itemize}
\item \emph{Exhaustive search} over all $2^{2^n}$ subfamilies of $\Pn([n])$ for
$n \le 4$, with a cardinality prune.
\item \emph{Branch and bound} for the maximum independent set of the $3$-uniform
hypergraph on the $2^n$ subsets whose hyperedges are the triples with
\emph{even} intersection, for $n = 5, 6, 7$. The bound greedily packs
vertex-disjoint forbidden triples (each permits at most $2$ of its three
vertices); vertices are ordered by conflict degree and the conflict structure is
precomputed from a parity table, so the $n = 7$ search runs in about ten seconds.
\end{itemize}

\begin{table}[h]
\centering
\begin{tabular}{r r r l}
\toprule
$n$ & true maximum $M(n)$ & claim $2^{n-3}$ & comparison \\
\midrule
$0$ & $1$ & --- (undefined) & --- \\
$1$ & $2$ & --- (undefined) & --- \\
$2$ & $2$ & --- (undefined) & --- \\
$3$ & $4$ & $1$  & claim too \textbf{small} by $3$ \\
$4$ & $5$ & $2$  & claim too \textbf{small} by $3$ \\
$5$ & $7$ & $4$  & claim too \textbf{small} by $3$ \\
$6$ & $8$ & $8$  & equal (coincidence only) \\
$7$ & $10$& $16$ & claim too \textbf{large} by $6$ \\
\bottomrule
\end{tabular}
\caption{True maxima $M(n)$ for $n = 0, \dots, 7$ versus the claim $2^{n-3}$.}
\end{table}

Equivalently,
\[
M(n) \;=\; 1,\, 2,\, 2,\, 4,\, 5,\, 7,\, 8,\, 10
\qquad (n = 0, 1, \dots, 7).
\]
The claim is too small at $n = 3, 4, 5$, too large at $n = 7$ (there is no
family of size $16$; the maximum is $10$), and agrees at $n = 6$ only by
coincidence. Hence $2^{n-3}$ is not the truth.

\section{The construction}

The extremal families for $n = 3, \dots, 7$ have the following shape. Fix a
point $x$ and let $Y = [n] \setminus \{x\}$. Put
\[
G \;=\; \{\emptyset\} \;\cup\; \bigl\{\{y\} : y \in Y\bigr\} \;\cup\;
\bigl\{\text{the edges of a matching on } Y\bigr\},
\qquad
F \;=\; \bigl\{\{x\} \cup g : g \in G\bigr\}.
\]

\begin{theorem}\label{thm:construction}
The family $F$ above is 3-wise odd-intersecting and has
\[
|F| \;=\; 1 + (n-1) + \Bigl\lfloor \frac{n-1}{2} \Bigr\rfloor
       \;=\; n + \Bigl\lfloor \frac{n-1}{2} \Bigr\rfloor .
\]
For $n = 3, 4, 5, 6, 7$ this equals $4, 5, 7, 8, 10$, i.e.\ $M(n)$.
\end{theorem}

\begin{proof}
Any three distinct members have the form $\{x\} \cup g_1$, $\{x\} \cup g_2$,
$\{x\} \cup g_3$ with distinct $g_i \in G$, and
\[
(\{x\}\cup g_1)\cap(\{x\}\cup g_2)\cap(\{x\}\cup g_3)
\;=\; \{x\} \cup (g_1 \cap g_2 \cap g_3).
\]
We claim $g_1 \cap g_2 \cap g_3 = \emptyset$. If some $g_i = \emptyset$ this is
clear. Otherwise each $g_i$ is a singleton $\{y\}$ or a matching edge. A
singleton $\{y\}$ equals at most one other $g_j$, and at most one matching edge
contains $y$; since $g_1, g_2, g_3$ are pairwise distinct, no point $y$ can
belong to all three. Two distinct matching edges are disjoint by definition of a
matching. Hence the intersection is empty and every triple meets in $\{x\}$, of
odd size $1$. The cardinality is $1 + |Y| + \lfloor |Y|/2 \rfloor$ with
$|Y| = n-1$.
\end{proof}

Thus the extremal families have a ``fixed point plus an even-intersecting family
on the remainder'' shape, not the affine-subspace shape the conjecture proposes.

\begin{remark}[Asymptotic lower bound]
Partitioning $Y$ into pairs and taking $G$ to be \emph{all} unions of these
pairs, the intersection $g_1 \cap g_2 \cap g_3$ is again a union of pairs (hence
of even size), so every triple still meets in $\{x\} \cup (\text{a union of
pairs})$, of odd size, and
\[
|F| \;=\; 2^{\lfloor (n-1)/2 \rfloor} \;\approx\; 2^{n/2}.
\]
This is an infinite family, so $M(n) \ge 2^{\lfloor (n-1)/2 \rfloor}$: the
maximum grows at least like $2^{n/2}$, an exponential of rate $n/2$ rather than
$n-3$. (The computed values $M(6) = 8 = 2^3$ and $M(7) = 10$ are consistent
with growth of that order. The lower bound alone does not determine $M(n)$ for
large $n$; what is decisive for the conjecture is the table of Section~4, which
refutes $2^{n-3}$ in both directions.)
\end{remark}

\section{The affine-subspace characterisation is false}

The phrase ``indicator functions of 3-dimensional affine subspaces'' admits two
readings. Identify a subset of $[n]$ with its indicator vector in $\F_2^n$, so
that a family of indicator vectors is a set of points of $\F_2^n$. Both readings
were tested by enumerating \emph{all} affine cosets of \emph{all} dimensions in
$\F_2^n$ for $n = 4, 5, 6, 7$.

\paragraph{Literal reading: a 3-dimensional affine subspace.}
Such a coset has $2^3 = 8$ members. Enumerating all $3$-dimensional subspaces
and all their cosets:

\begin{center}
\begin{tabular}{r r r r}
\toprule
$n$ & $3$-dim.\ subspaces & cosets & cosets that are 3-wise odd \\
\midrule
$4$ & $15$   & $30$    & $0$ \\
$5$ & $155$  & $620$   & $0$ \\
$6$ & $1395$ & $11160$ & $0$ \\
$7$ & $11811$& $188976$& $105$ \\
\bottomrule
\end{tabular}
\end{center}

For $n = 4, 5, 6$ there is \emph{no} 3-dimensional affine coset that is
3-wise odd-intersecting. For $n = 7$ there are $105$ such cosets, but every one
has size $8 < M(7) = 10$, so none is extremal.

\paragraph{Codimension-3 reading: a coset of size $2^{n-3}$.}
Here $n = 4$: cosets of size $2$, all vacuously 3-wise odd, but $2 < M(4) = 5$;
$n = 5$: cosets of size $4$, of which $460$ are 3-wise odd, but
$4 < M(5) = 7$; $n = 6$: cosets of size $8$, and \emph{none} is 3-wise odd;
$n = 7$: such a coset would have size $16 > M(7) = 10$, impossible.

\begin{corollary}
No affine coset of any dimension is an extremal family for $n = 4, 5, 6, 7$. If
$n \in \{4,5,7\}$ then $M(n) = 5, 7, 10$ is not a power of two, while every
affine coset has power-of-two size, so no coset can even have $M(n)$ members. At
$n = 6$, where $M(6) = 8 = 2^3$ is a power of two, the only cosets of size $8$
are $3$-dimensional, and none of them is 3-wise odd-intersecting.
\end{corollary}

\begin{remark}[Ambiguity note]
``Indicator functions of 3-dimensional affine subspaces'' is garbled: a
3-dimensional affine subspace has $2^3 = 8$ points, whereas $2^{n-3}$ is the
size of a codimension-3 coset. Each indicator function of a 3-dimensional
affine subspace is a \emph{member} of the proposed extremal family, while
$2^{n-3}$ is a \emph{number of members}; the two readings are different
statements, and neither yields valid extremal families.
\end{remark}

\begin{remark}
The $n = 4$ witness is itself an affine coset, but of dimension $2$: its
indicator vectors satisfy
$F = \{1100, 1010, 1001, 1111\} = 1100 + \{0000, 0110, 0101, 0011\}$,
a coset of the span of $0110$ and $0101$. Its size is $2^2 = 4$, its dimension
is $2$ (not $3$), and it is not extremal because $M(4) = 5$. So being an affine
coset is neither sufficient nor, at the claimed dimension, possible.
\end{remark}

\section{Formalisation}

The $n = 4$ witness is formalised in the accompanying Lean~4 project under
\texttt{lean4/}, using core Lean only (\texttt{import Std}, no Mathlib, no
\texttt{sorry}). Members are represented as indicator vectors of type
\texttt{List Bool}. The main statements are

\begin{quote}\ttfamily
def interSize (a b c : List Bool) : Nat\\
def TripleOdd (F : List (List Bool)) : Prop\\
def F : List (List Bool) :=\\
\hspace*{2em}[[true,true,false,false], [true,false,true,false],\\
\hspace*{2em} [true,false,false,true], [true,true,true,true]]\\[2pt]
theorem witness\_oddTriples : TripleOdd F\\
theorem triple\_123 : interSize m12 m13 m14 = 1\\
theorem triple\_124 : interSize m12 m13 m1234 = 1\\
theorem triple\_134 : interSize m12 m14 m1234 = 1\\
theorem triple\_234 : interSize m13 m14 m1234 = 1\\
theorem witness\_length : F.length = 4\\
theorem witness\_breaks\_bound : F.length > 2 \^{} (4 - 3 : Nat)\\
theorem witness\_check : checkFamily F = true\\
theorem n4\_max\_at\_least\_five :\\
\hspace*{2em}TripleOdd G5 \(\wedge\) G5.length = 5 \(\wedge\) G5.length > 2 \^{} (4 - 3 : Nat)
\end{quote}

\noindent Here \texttt{TripleOdd} is the conjecture's property, \texttt{F} is
the four-member witness, and \texttt{G5} is the five-member family of
Corollary~3.1. The four triple-intersection theorems are the explicit check
that the witness meets the property; \texttt{witness\_breaks\_bound} is the
inequality $4 > 2$. Building and auditing:

\begin{quote}\ttfamily
\$ cd lean4\\
\$ lake build\\
\$ lake env lean Check.lean
\end{quote}

\noindent The audit reports no \texttt{sorryAx}; the concrete computational
theorems depend on no axioms at all, and the quantified ones use only
\texttt{propext} and \texttt{Quot.sound}. The exact maxima $M(n)$ and the
affine-coset enumeration are finite computations carried out by
\texttt{reproduce.py}; they are not needed for the refutation, which is already
complete at $n = 4$. See \texttt{lean4/README.md} for the full statement table
and scope.

\section{Reproducibility}

The checks are carried out by \texttt{reproduce.py}, which uses the Python~3
standard library only:

\begin{quote}\ttfamily
\$ python3 reproduce.py\\
... \textit{(the witness, the maxima, the constructions, the affine cosets)}\\
PASS: all checks verified; conjecture 00000002175 is FALSE.
\end{quote}

It verifies the witness's four triple intersections (each equal to $\{1\}$),
computes $M(0..4)$ exhaustively and $M(5..7)$ by branch and bound (cross-checking
the two methods at $n = 4$), checks the constructions of Section~5, and tests
every affine coset for $n = 4, \dots, 7$. It prints \texttt{PASS} and exits $0$
exactly when every check holds. The document itself is built with
\texttt{tectonic --outdir build main.tex}.

\end{document}
